%% ===== 3 Related Work =====
\section{Related Work}
\label{sec:related}

\paragraph{\new{Closest related studies.}} Chen et al.\ present the most recent \new{systematic literature review (SLR)} on neural code translation, analyzing 57 primary studies from seven perspectives, including evaluation subjects and metrics~\cite{nctslr}. Their \new{catalog} records which metrics appear; we take evaluation itself as the object of analysis, classifying where oracles come from and how reliable they are, and contrasting settings with and without tests. Jiao et al.\ propose an evaluation taxonomy and benchmark for neural code translation~\cite{jiao23}; their taxonomy classifies translation task types in the pre-LLM era, whereas we map evaluation practice in the LLM and agent era, including LLM-generated tests and LLM judges. A security-focused \new{systematization of knowledge (SoK)} of AI4Code sketches verification techniques for translation within one subsection~\cite{ai4codesok}; we give the topic a dedicated, corpus-based treatment. A recent systematic mapping study of generative AI for code translation maps the publication landscape by year, venue, research type, dataset and model~\cite{genaisms}; it does not analyze oracle origin or reliability, which are our central axes.

\paragraph{The oracle problem.} Barr et al.\ survey the test oracle problem in general~\cite{barr}. Migration is a special case in which the original program is available as a reference, so differential and equivalence-based oracles become possible; our map shows how rarely that advantage is used \new{(13 of 49 papers; Figure~\ref{fig:oracle-freq})}.

\paragraph{\new{Individual studies.}} Individual benchmarks and techniques already work around the oracle problem: coverage gates against reward hacking~\cite{freshbrew}, tests translated with the code~\cite{alphatrans}, differential fuzzing without predefined tests~\cite{dristi}, and benchmarks for equivalence judgments~\cite{polyhuman}. \new{These works document the defenses and failure modes that we consolidate. Two of them, FreshBrew and AlphaTrans, are in our corpus. Of the other two, one studies refactoring within a single language and fails our criteria, and the other appeared at ESEM, outside our venues.}
